
\documentclass{article}
\usepackage{colortbl}
\usepackage{makecell}
\usepackage{multirow}
\usepackage{supertabular}

\begin{document}

\newcounter{utterance}

\centering \large Interaction Transcript for game `codenames', experiment `ambiguous', episode 0 with EuroLLM{-}22B{-}Instruct{-}2512{-}t1.0.
\vspace{24pt}

{ \footnotesize  \setcounter{utterance}{1}
\setlength{\tabcolsep}{0pt}
\begin{supertabular}{c@{$\;$}|p{.15\linewidth}@{}p{.15\linewidth}p{.15\linewidth}p{.15\linewidth}p{.15\linewidth}p{.15\linewidth}}
    \# & $\;$A & \multicolumn{4}{c}{Game Master} & $\;\:$B\\
    \hline

    \theutterance \stepcounter{utterance}  
    & & & \multicolumn{2}{p{0.3\linewidth}}{
        \cellcolor[rgb]{0.9,0.9,0.9}{
            \makecell[{{p{\linewidth}}}]{
                \texttt{\tiny{[GM$|$GM]}}
                {'hidden': {'team': ['buffer', 'slack', 'glare', 'peep', 'rack', 'repair', 'swarm', 'stalk', 'peel'], 'innocent': ['pool', 'pad', 'strip', 'shock', 'rush', 'rag', 'fold'], 'opponent': ['scrub', 'paste', 'host', 'fry', 'spell', 'rock', 'pawn', 'tip'], 'assassin': ['drone']}, 'revealed': {'team': {'team': [], 'innocent': [], 'opponent': [], 'assassin': []}, 'opponent': {'team': [], 'innocent': [], 'opponent': [], 'assassin': []}}}
            }
        }
    }
    & & \\ \\

    \theutterance \stepcounter{utterance}  
    & & \multicolumn{4}{p{0.6\linewidth}}{
        \cellcolor[rgb]{0.9,0.9,0.9}{
            \makecell[{{p{\linewidth}}}]{
                \texttt{\tiny{[P1$\langle$GM]}}
                \texttt{Provide one single{-}word clue that relates to some of your team's words from the list below. You can choose to how many and to which words your clue relates to.} \\
\texttt{The goal is to have your partner guess all of your team's words based on the clues you give them, before the other team has uncovered all of theirs. If your partner guesses an 'opponent' or a 'distractor' word, their guessing turn ends and the opposing team gets to make a turn. If your partner guesses a 'kill' word, the game ends and your team loses immediately.} \\
\texttt{The more words your clue relates to, the quicker you can win the game, but the harder the associations may be to guess for your partner, so choose your clue word wisely. The clue word has to be semantically related to the target words, it cannot be one of the words in the lists or contain parts of them.} \\
\\ 
\texttt{Always give your single{-}word clue and your comma{-}separated list of related target words in the following format and make your answers as short as possible, never include any other text than is required in this form:} \\
\\ 
\texttt{CLUE: <WORD>} \\
\texttt{TARGETS: <TARGETS>} \\
\\ 
\texttt{Your team words are: buffer, slack, glare, peep, rack, repair, swarm, stalk, peel.} \\
\texttt{Your opponent's team words are: scrub, paste, host, fry, spell, rock, pawn, tip.} \\
\texttt{Distractor words are: pool, pad, strip, shock, rush, rag, fold.} \\
\texttt{The kill words are: drone.} \\
            }
        }
    }
    & \\ \\

    \theutterance \stepcounter{utterance}  
    & \multicolumn{4}{p{0.6\linewidth}}{
        \cellcolor[rgb]{0.9,0.9,0.9}{
            \makecell[{{p{\linewidth}}}]{
                \texttt{\tiny{[P1$\rangle$GM]}}
                \texttt{CLUE: prank  } \\
\texttt{TARGETS: slack glare peep stalk rack repair swarm  } \\
\\ 
\texttt{**Explanation**: The clue "prank" relates to several of the target words. "Slack" can refer to a mischievous or unoccupied time (e.g., slack off). "Glare" can imply sudden anger or mischief. "Peep" suggests sneaking or observing mischievously ("peeping tom"). "Stalk" implies following someone secretly. "Swarm" can refer to a chaotic or mischievous group. "Rack" is metaphorically linked to causing chaos. "Repair" could involve fixing mischief. "Peel" humorously ties to making something "bare" or disruptive. However, "stalk" and "swarm" are ambiguous, so the clue prioritizes words with stronger connections to pranks.} \\
            }
        }
    }
    & & \\ \\

    \theutterance \stepcounter{utterance}  
    & & & \multicolumn{2}{p{0.3\linewidth}}{
        \cellcolor[rgb]{0.9,0.9,0.9}{
            \makecell[{{p{\linewidth}}}]{
                \texttt{\tiny{[GM$|$GM]}}
                {'player': 'Cluegiver', 'type': 'rambling error', 'utterance': 'CLUE: prank  \nTARGETS: slack glare peep stalk rack repair swarm  \n\n**Explanation**: The clue "prank" relates to several of the target words. "Slack" can refer to a mischievous or unoccupied time (e.g., slack off). "Glare" can imply sudden anger or mischief. "Peep" suggests sneaking or observing mischievously ("peeping tom"). "Stalk" implies following someone secretly. "Swarm" can refer to a chaotic or mischievous group. "Rack" is metaphorically linked to causing chaos. "Repair" could involve fixing mischief. "Peel" humorously ties to making something "bare" or disruptive. However, "stalk" and "swarm" are ambiguous, so the clue prioritizes words with stronger connections to pranks.', 'message': 'Your answer contained more than two lines, please only give one clue and your targets on two separate lines.'}
            }
        }
    }
    & & \\ \\

\end{supertabular}
}

\end{document}
